% ===========================================================================
%  Personal shorthand macros
% ---------------------------------------------------------------------------
%  Number sets, common upright symbols and operators, plus \permu for drawing
%  permutation diagrams:  \permu{1,2,3}{2,3,1}  (top row, images).
% ===========================================================================

% ---------------------------------------------------------------------------
%  Locating a file from any directory
% ---------------------------------------------------------------------------
%  \projectfindfile{path} puts the first of
%      path      ../path      ../../path
%  that actually exists into \projectfoundfile, or leaves it empty.
%
%  Needed because \IfFileExists also searches \input@path - which exercise
%  sheets and slides extend with "../" so they can load style/project - while
%  \includegraphics searches \graphicspath. The two disagree for documents
%  that do not sit in the project root, which silently loses the image.
%  Resolving the path once, here, keeps them in step.
% ---------------------------------------------------------------------------
\newcommand{\projectfoundfile}{}

\newcommand{\projectfindfile}[1]{%
    \begingroup
        \let\input@path\@undefined
        \IfFileExists{#1}%
            {\gdef\projectfoundfile{#1}}%
            {\IfFileExists{../#1}%
                {\gdef\projectfoundfile{../#1}}%
                {\IfFileExists{../../#1}%
                    {\gdef\projectfoundfile{../../#1}}%
                    {\gdef\projectfoundfile{}}}}%
    \endgroup
}

% Include \projectlogo if it can be found, otherwise leave a placeholder frame
% of the same width. #1 = \includegraphics options.
\newcommand{\projectlogobox}[1]{%
    \projectfindfile{\projectlogo}%
    \ifx\projectfoundfile\empty
        \fbox{\rule{0pt}{3cm}\rule{0.28\textwidth}{0pt}}%
    \else
        \includegraphics[#1]{\projectfoundfile}%
    \fi
}

% Centred logo, or nothing at all if there is none. Everything stays inside
% this one macro: a conditional wrapped around \begin{center} would break
% inside a beamer frame, whose body is collected and re-scanned.
\newcommand{\projectlogocenter}[1]{%
    \projectfindfile{\projectlogo}%
    \ifx\projectfoundfile\empty\else
        {\par\centering\includegraphics[#1]{\projectfoundfile}\par}%
    \fi
}

\newcommand{\N}{\mathbb{N}}
\newcommand{\Z}{\mathbb{Z}}
\newcommand{\Q}{\mathbb{Q}}
\newcommand{\R}{\mathbb{R}}
\newcommand{\C}{\mathbb{C}}
\newcommand{\D}{\mathrm{d}}
\newcommand{\E}{\mathrm{e}}
\newcommand{\I}{\mathrm{i}}
\newcommand{\bs}{\boldsymbol}
\DeclareMathOperator{\diag}{diag}
\DeclareMathOperator{\lcm}{lcm}
\newcommand*\laplaceop{\mathop{}\!\mathbin\bigtriangleup}
\newcommand*\quabla{\mathop{}\!\mathbin\Box}

\ExplSyntaxOn

\msg_new:nnn { permu } { length-mismatch }
  { The~lists~in~\string\permu~must~have~the~same~length. }

\msg_new:nnn { permu } { missing-image }
  { The~image~'#1'~in~\string\permu~does~not~occur~in~the~top~row. }

\seq_new:N \l__permu_top_seq
\seq_new:N \l__permu_img_seq
\int_new:N \l__permu_len_int
\int_new:N \l__permu_target_int

\cs_new_protected:Npn \__permu_find_index:n #1
{
    \int_zero:N \l__permu_target_int
    \seq_map_indexed_inline:Nn \l__permu_top_seq
    {
        \str_if_eq:eeT { \tl_trim_spaces:n { ##2 } } { \tl_trim_spaces:n { #1 } }
        {
            \int_set:Nn \l__permu_target_int { ##1 }
        }
    }
}

\cs_new_protected:Npn \__permu_draw:nn #1 #2
{
    \seq_set_split:Nnn \l__permu_top_seq { , } { #1 }
    \seq_set_split:Nnn \l__permu_img_seq { , } { #2 }
    \int_set:Nn \l__permu_len_int { \seq_count:N \l__permu_top_seq }

    \int_compare:nNnF
        { \l__permu_len_int }
        =
        { \seq_count:N \l__permu_img_seq }
        { \msg_error:nn { permu } { length-mismatch } }

    \begin{tikzpicture}[baseline={(0,0)}]
        \seq_map_indexed_inline:Nn \l__permu_top_seq
        {
            \node (permutop##1) at ({1.15*(##1-1)},0.65) {$\scriptstyle ##2$};
            \node (permubot##1) at ({1.15*(##1-1)},-0.65) {$\scriptstyle ##2$};
        }

        \seq_map_indexed_inline:Nn \l__permu_img_seq
        {
            \__permu_find_index:n { ##2 }

            \int_compare:nNnTF { \l__permu_target_int } > { 0 }
            {
                {
                    \pgfsetlinewidth{0.9pt}
                    \draw
                        (permutop##1.south) -- (permubot\int_use:N \l__permu_target_int.north);
                }
            }
            {
                \msg_error:nnn { permu } { missing-image } { ##2 }
            }
        }
    \end{tikzpicture}
}

\NewDocumentCommand{\permu}{m m}
{
    \__permu_draw:nn { #1 } { #2 }
}

\ExplSyntaxOff
